\begin{lemma}
\label{lemma-unbounded-right-derived}
Let $F : \mathcal{A} \to \mathcal{B}$ be a left exact functor of
abelian categories. Assume
\begin{enumerate}
\item every object of $\mathcal{A}$ is a subobject of an object
which is right acyclic for $F$,
\item there exists an integer $n \geq 0$ such that $R^nF = 0$,
\end{enumerate}
Then
\begin{enumerate}
\item $RF : D(\mathcal{A}) \to D(\mathcal{B})$ exists,
\item any complex consisting of right acyclic objects for $F$ computes $RF$,
\item any complex is the source of a quasi-isomorphism into a complex
consisting of right acyclic objects for $F$,
\item for $E \in D(\mathcal{A})$
\begin{enumerate}
\item $H^i(RF(\tau_{\leq a}E) \to H^i(RF(E))$ is an isomorphism
for $i \leq a$,
\item $H^i(RF(E)) \to H^i(RF(\tau_{\geq b - n + 1}E))$ is an isomorphism
for $i \geq b$,
\item if $H^i(E) = 0$ for $i \not \in [a, b]$ for some
$-\infty \leq a \leq b \leq \infty$, then $H^i(RF(E)) = 0$
for $i \not \in [a, b + n - 1]$.
\end{enumerate}
\end{enumerate}
\end{lemma}

\begin{proof}
Note that the first assumption implies that
$RF : D^+(\mathcal{A}) \to D^+(\mathcal{B})$ exists, see
Proposition \ref{proposition-enough-acyclics}.
Let $A$ be an object of $\mathcal{A}$. Choose an injection $A \to A'$
with $A'$ acyclic. Then we see that $R^{n + 1}F(A) = R^nF(A'/A) = 0$ by
the long exact cohomology sequence. Hence we conclude that $R^{n + 1}F = 0$.
Continuing like this using induction we find that $R^mF = 0$ for all
$m \geq n$.

\medskip\noindent
We are going to use Lemma \ref{lemma-replace-resolution} with the function
$d : \Ob(\mathcal{A}) \to \{0, 1, 2, \ldots \}$ given by
$d(A) = \max \{0\} \cup \{i \mid R^iF(A) \not = 0\}$.
The first assumption of Lemma \ref{lemma-replace-resolution}
is our assumption (1). The second assumption of
Lemma \ref{lemma-replace-resolution} follows from the fact
that $RF(A \oplus B) = RF(A) \oplus RF(B)$. The third assumption of
Lemma \ref{lemma-replace-resolution} follows from the long exact
cohomology sequence. Hence for every complex $K^\bullet$ there exists a
quasi-isomorphism $K^\bullet \to L^\bullet$ into a complex of
objects right acyclic for $F$. This proves statement (3).

\medskip\noindent
We claim that if $L^\bullet \to M^\bullet$ is a quasi-isomorphism of
complexes of right acyclic objects for $F$, then
$F(L^\bullet) \to F(M^\bullet)$
is a quasi-isomorphism. If we prove this claim then we get statements
(1) and (2) of the lemma by
Lemma \ref{lemma-find-existence-computes}.
To prove the claim pick an integer $i \in \mathbf{Z}$.
Consider the distinguished triangle
$$
\sigma_{\geq i - n - 1}L^\bullet \to
\sigma_{\geq i - n - 1}M^\bullet \to Q^\bullet,
$$
i.e., let $Q^\bullet$ be the cone of the first map.
Note that $Q^\bullet$ is bounded below and that
$H^j(Q^\bullet)$ is zero except possibly for $j = i - n - 1$
or $j = i - n - 2$. We may apply $RF$ to $Q^\bullet$.
Using the second spectral sequence of
Lemma \ref{lemma-two-ss-complex-functor}
and the assumed vanishing of cohomology (2) we conclude
that $H^j(RF(Q^\bullet))$ is zero except possibly for
$j \in \{i - n - 2, \ldots, i - 1\}$. Hence we see that
$RF(\sigma_{\geq i - n - 1}L^\bullet) \to RF(\sigma_{\geq i - n - 1}M^\bullet)$
induces an isomorphism of cohomology objects in degrees $\geq i$.
By Proposition \ref{proposition-enough-acyclics} we know that
$RF(\sigma_{\geq i - n - 1}L^\bullet) = \sigma_{\geq i - n - 1}F(L^\bullet)$
and
$RF(\sigma_{\geq i - n - 1}M^\bullet) = \sigma_{\geq i - n - 1}F(M^\bullet)$.
We conclude that $F(L^\bullet) \to F(M^\bullet)$
is an isomorphism in degree $i$ as desired.

\medskip\noindent
Part (4)(a) follows from Lemma \ref{lemma-negative-vanishing}.

\medskip\noindent
For part (4)(b) let $E$ be represented by the complex $L^\bullet$
of objects right acyclic for $F$. By part (2) $RF(E)$ is represented
by the complex $F(L^\bullet)$ and $RF(\sigma_{\geq c}L^\bullet)$
is represented by $\sigma_{\geq c}F(L^\bullet)$. Consider the
distinguished triangle
$$
H^{b - n}(L^\bullet)[n - b] \to
\tau_{\geq b - n}L^\bullet \to
\tau_{\geq b - n + 1}L^\bullet
$$
of Remark \ref{remark-truncation-distinguished-triangle}.
The vanishing established above gives that
$H^i(RF(\tau_{\geq b - n}L^\bullet))$ agrees with
$H^i(RF(\tau_{\geq b - n + 1}L^\bullet))$ for $i \geq b$.
Consider the short exact sequence of complexes
$$
0 \to
\Im(L^{b - n - 1} \to L^{b - n})[n - b] \to
\sigma_{\geq b - n}L^\bullet \to
\tau_{\geq b - n}L^\bullet \to 0
$$
Using the distinguished triangle associated to this
(see Section \ref{section-canonical-delta-functor})
and the vanishing as before we conclude that
$H^i(RF(\tau_{\geq b - n}L^\bullet))$ agrees with
$H^i(RF(\sigma_{\geq b - n}L^\bullet))$ for $i \geq b$.
Since the map $RF(\sigma_{\geq b - n}L^\bullet) \to RF(L^\bullet)$
is represented by $\sigma_{\geq b - n}F(L^\bullet) \to F(L^\bullet)$
we conclude that this in turn agrees with $H^i(RF(L^\bullet))$
for $i \geq b$ as desired.

\medskip\noindent
Proof of (4)(c). Under the assumption on $E$ we have
$\tau_{\leq a - 1}E = 0$ and we get the vanishing
of $H^i(RF(E))$ for $i \leq a - 1$ from part (4)(a).
Similarly, we have $\tau_{\geq b + 1}E = 0$ and hence
we get the vanishing of $H^i(RF(E))$ for $i \geq b + n$ from
part (4)(b).
\end{proof}